% SOURCES: docs/0.2.0v/plan.md (§10 Tasks, §17 Milestones, §9 Library-to-Feature Map);
%          docs/CHANGELOG.md; docs/DECISIONS.md; docs/1.0.0v/supervisor-report.md (§1)

\chapter{Project Implementation Plan}
\label{ch:plan}

\section{Milestone Breakdown}
The work was divided into five gated milestones plus the v1.0.0 demo cut
(Table~\ref{tab:milestones}). The ordering encodes the technical dependencies of
the system: the Document Model and the store had to exist before anything could be
rendered or generated (M1); themable output required the token system (M2);
composition and responsiveness required presets and per-breakpoint style (M3); the
runtime, SEO, and the accessibility gate hardened the output (M4); and polish plus
a released artifact completed the sprint (M5). Each milestone was ``done'' only
when its gating tasks were complete and all quality gates were green.

\begin{table}[H]
  \centering
  \caption{Milestone plan and definition of done.}
  \label{tab:milestones}
  \begin{tabularx}{\linewidth}{@{}l l X@{}}
    \toprule
    \textbf{Milestone} & \textbf{Theme} & \textbf{Definition of done (summary)} \\
    \midrule
    M1 & Foundation            & Types + store + recursive canvas + undo/redo + persistence + basic generator + export shell. \\
    M2 & Tokens + Themes       & Token registry + theme toggle + auto-layout + \texttt{:root} block + \texttt{var()} refs. \\
    M3 & Composition + Resp.   & 6+ presets + hover/focus states + breakpoint switcher + multi-select + image upload. \\
    M4 & Runtime + Hardening   & Full runtime snippets + SEO + validation console + \texttt{axe-core} gate + templates. \\
    M5 & Polish + Demo         & Polish + performance budgets + release artifact + rehearsed demo. \\
    v1.0.0 & Demo cut          & draw-to-create + match-layout + headless MCP server; released with native installers. \\
    \bottomrule
  \end{tabularx}
\end{table}

\section{Task Allocation and the Critical Path}
% SOURCES: docs/0.2.0v/plan.md (§10); CLAUDE.md (Module Ownership)
Tasks were identified by owner and area --- for example \texttt{I-DOC-*} for the
Document Model, \texttt{I-GEN-*} for the generator, \texttt{I-RUN-*} for runtime
snippets, \texttt{Y-*} for the store, and \texttt{L-*} for the user interface ---
each with a priority and explicit dependency links. The critical path ran through
the Document Model: because the types (C1), schemas (C2), and operations (C3) are
consumed by every other lane, the \texttt{I-DOC} tasks were front-loaded and had to
land before the store, canvas, and generator could be written against a stable
shape. Table~\ref{tab:allocation} summarises the allocation of areas to owners.

\begin{table}[H]
  \centering
  \caption{Allocation of work areas to owners.}
  \label{tab:allocation}
  \small
  \begin{tabularx}{\linewidth}{@{}l X@{}}
    \toprule
    \textbf{Owner} & \textbf{Areas (task prefixes)} \\
    \midrule
    Ibrahim Haski      & Document Model (\texttt{I-DOC}), Electron shell (\texttt{I-ELE}), generator (\texttt{I-GEN}), runtime (\texttt{I-RUN}), SEO (\texttt{I-SEO}), export (\texttt{I-EXP}), templates (\texttt{I-TPL}), build/CI (\texttt{I-BLD}). \\
    Yousef Deeb        & Zustand stores, history via \texttt{immer} patches, persistence, autosave/crash recovery, external-change reload (\texttt{Y-*}); draw / match / MCP surfaces (v1.0.0). \\
    Abd Alrazak Qahwaji & Canvas, properties/layers/tokens panels, sidebar, topbar, live validation console, semantic-inference hints (\texttt{L-*}). \\
    \bottomrule
  \end{tabularx}
\end{table}

\section{Deviations from the Plan}
\label{sec:deviations}
% SOURCES: DECISIONS.md (D-01..D-04)
A plan is only credible if its deviations are recorded honestly. Four decisions
in the project's append-only decision log changed the plan, and each is
summarised here.

\begin{itemize}[leftmargin=1.5em]
  \item \textbf{D-01 --- Code signing cancelled permanently.} Windows Authenticode
        signing and macOS notarization (task \texttt{I-BLD-05}) were cut from scope,
        not merely deferred: the team has no code-signing certificate and no Apple
        Developer ID and decided the cost was not justified. Builds therefore ship
        \emph{unsigned}, and the README documents how users bypass the resulting
        one-time operating-system trust prompt.
  \item \textbf{D-02 / D-04 --- Releases shipped ahead of the manual demo
        rehearsal and the render-layer performance sign-off.} The v0.3.0 and
        v1.0.0 tags were cut while two \emph{manual} items remained open: the
        rehearsed demo on the packaged build, and the empirical render-layer
        performance capture (drag FPS, theme-toggle latency, and similar). These
        are not code defects --- the automated suite is green --- but a conscious,
        documented waiver, tracked as the top post-release item (issue~\#95).
  \item \textbf{D-03 --- Three Tier-2 authoring surfaces brought forward into
        v1.0.0.} Draw-to-create, match-layout, and the MCP server were originally
        Tier-2 (out of scope), but were merged and shipped as the headline of the
        demo cut because they are the most demonstrable features in the product
        and, being pure and headless at their core and reusing the existing
        operations, added capability without adding architectural risk.
\end{itemize}

Recording these deviations, with their rationale and accepted consequences, is
part of the methodology (Chapter~\ref{ch:methodology}): the decision log is the
mechanism that keeps ``what we planned'' and ``what we shipped'' reconcilable.
